We develop a mathematical framework for analyzing chain-of-thought (CoT) reasoning
as a discrete dynamical system on the probability simplex equipped with the
Wasserstein metric. We model CoT as a Markov transition operator $T$ acting on
probability distributions over a finite state space of reasoning states, and
establish convergence rate bounds in two regimes. In the contractive regime, we
prove that when the Dobrushin coefficient $\delta(T)$ is strictly less than one,
the iterates converge exponentially to the unique stationary distribution in the
$W_1$ metric, with rate governed by a finer Wasserstein contraction coefficient.
In the metastable regime, for transition matrices with $(K,\varepsilon)$-metastable
cluster structure, we derive an exact formula for the spectral gap
$\gamma = K\varepsilon/(K-1) + O(\varepsilon^2)$ using perturbation theory on
the aggregated inter-cluster chain, yielding precise mixing time asymptotics.
We characterize a phase transition in convergence behavior at a critical coupling
$\varepsilon_c$ separating slow and fast mixing regimes. Our proofs combine
Kantorovich--Rubinstein duality, the Banach fixed-point theorem, and matrix
perturbation theory. All theoretical predictions are verified computationally:
the correlation between predicted spectral convergence rates and empirical $W_2$
decay rates is $r = 0.979$ across 30 random transition matrices, and the
predicted spectral gaps for metastable chains match numerical eigenvalue
computations to machine precision.
